There exists a compact set of transfinite diameter one, outside every closed unit disk, which allows $n$ components for every positive integer $n$. On the other hand, for each fixed nonempty compact set $F$ of capacity less than one there are constants $c(F)>0$ and $N(F)$ such that every monic degree-$n$ polynomial with roots in $F$ has at most $(1-c(F))n$ components once $n\ge N(F)$.

The second assertion is asymptotic: for $n=1$ there is always one component, so a bound $(1-c)n$ with $c>0$ cannot hold at every degree. This is exactly the quantifier in Theorem 2.1(a) of Ghosh--Ramachandran, \emph{On the number of components of polynomial lemniscates}, \url{https://arxiv.org/abs/2312.13673}. We reconstruct its compactness argument, and give an explicit example for the first assertion. We do not assert that every capacity-one set has the first property, or determine the optimal constant solely from the capacity.
